% Зургийн өнгөний нэг палитр (2026.11.20). Багш нар 11/17-нд «графын өнгө бүгд ижил байна» гэсэн:
% өмнөх бүдэг хольц (blue!10, violet!13 г.м.) дэлгэц дээр ялгардаггүй байв. Одоо бүлэг бүр өөр өнгөтэй:
% бүдэг дүүргэлт = үйлдэл, алхам, мэдээлэл; тод дүүргэлт = оролцогч, гадны үйлчилгээ.
% Хар бичгийн тодрол бүх дүүргэлт дээр 4.5:1-ээс их; хүрээний өнгө цагаан дээр 3:1-ээс их.
% Өнгө дангаараа утга илэрхийлэхгүй — зураг бүр бичиг, тайлбартай.
% Энэ файлыг main.tex ба tools/export_thesis_figure.py хоёулаа уншина.
\definecolor{nmBlue}{HTML}{BBDEFB}    % системийн алхам; зочны хэрэглээний тохиолдол; вэб сайтаас авсан мэдээлэл
\definecolor{nmBlueS}{HTML}{64B5F6}   % тодруулсан хөх (задалж үзүүлсэн хүрээ)
\definecolor{nmBlueD}{HTML}{0D47A1}
\definecolor{nmTeal}{HTML}{80CBC4}    % хэрэглэгч, зочин
\definecolor{nmTealD}{HTML}{00695C}
\definecolor{nmOrange}{HTML}{FFB74D}  % гадны үйлчилгээ
\definecolor{nmOrangeL}{HTML}{FFE0B2} % гадны үйлчилгээгээр гүйцэтгэх үйлдэл
\definecolor{nmOrangeD}{HTML}{BF360C}
\definecolor{nmGreen}{HTML}{C8E6C9}   % мэдээлэл; сургуулийн зүгээс өгсөн мэдээлэл
\definecolor{nmGreenS}{HTML}{81C784}  % мэдээллийг оруулдаг сургуулийн ажилтан; тодруулсан ногоон
\definecolor{nmGreenD}{HTML}{1B5E20}
\definecolor{nmViolet}{HTML}{E1BEE7}  % хэнд, хаана хандахыг заах; модны үндэс
\definecolor{nmVioletD}{HTML}{6A1B9A}
